% id: 108-13
% book: 베이직쎈 확률과 통계 · 06 이항분포와 정규분포
% section: 개념 31 정규분포곡선의 성질
% passage: 오른쪽 그림에서 세 곡선 $\mathrm{A}$, $\mathrm{B}$, $\mathrm{C}$는 각각 정규분포 $\mathrm{N}(m_{\mathrm{A}}{,}\mkern3mu\mathopen{}{\sigma_{\mathrm{A}}}^2)$, $\mathrm{N}(m_{\mathrm{B}}{,}\mkern3mu\mathopen{}{\sigma_{\mathrm{B}}}^2)$, $\mathrm{N}(m_{\mathrm{C}}{,}\mkern3mu\mathopen{}{\sigma_{\mathrm{C}}}^2)$을 따르는 세 확률변수의 정규분포곡선이다. 두 곡선 $\mathrm{A}$, $\mathrm{B}$의 모양이 같고, 두 곡선 $\mathrm{B}$, $\mathrm{C}$의 대칭축이 같을 때, 다음 두 수의 대소를 비교하시오.
% figure: tikz:fig-108-11
% answer: 
% answer_source: 
% variant_level: 0 (원본 전사)
% 이 파일은 build.mjs 가 items.json 에서 만든다. 고칠 때는 items.json 을 고친다.
\smash{\raisebox{1.5mm}{\dmkichul{베이직쎈 확률과 통계 108쪽 13번}}}
$\sigma_{\mathrm{A}}\ \blank\ \sigma_{\mathrm{B}}$
